\clearpage\fancyhead[L]{\small KOMET.}\hypertarget{native-407-part-lower}{}\label{native-407-part-lower}
\begin{center}\large KOMET.\end{center}
\begin{center}\hrule\vspace{4pt}\fontsize{10}{12}\selectfont V. \textsc{Lindenau} u. \textsc{Bohnenberger} Zeitschrift für Astronomie u. verwandte Wiss.\par Band V. Seite 276..277. März und April 1818.\par\vspace{4pt}\hrule\vspace{4pt}\end{center}
\par\indent Aus den Beobachtungen des Kometen von Dr. \textsc{Olbers} vom 3. bis 28. März habe ich in diesen Tagen die folgenden parabolischen Elemente abgeleitet. Sie weichen nur wenig von denen ab, die \textsc{Olbers} selbst aus seinen Beobachtungen vom 3. und 13. März und der mir noch unbekannten Beobachtung aus Marseille vom 18. Januar erhalten hat.

\begin{center}\fontsize{10}{12}\selectfont
Durchgang durch das Perihel am 26,8991. Februar 1818, Göttinger Zeit.
\par\smallskip\begin{tabular}{lrr}
Logarithmus des Abstands im Perihel & $\cdot\;\cdot\;\cdot\;\cdot\;\cdot\;\cdot\;\cdot\;\cdot$ & $0.07937$\\
Neigung der Bahn & $\cdot\;\cdot\;\cdot\;\cdot\;\cdot\;\cdot\;\cdot\;\cdot\;\cdot\;\cdot\;\cdot\;\cdot\;\cdot\;\cdot\;\cdot$ & $90^\circ\;0'\;0''$\\
Aufsteigender Knoten & $\cdot\;\cdot\;\cdot\;\cdot\;\cdot\;\cdot\;\cdot\;\cdot\;\cdot\;\cdot\;\cdot\;\cdot\;\cdot\;\cdot$ & $70\;5\;12$\\
\end{tabular}
\par\smallskip Länge des Perihels, die Bewegung als rechtläufig betrachtet $183\;22\;58$
\end{center}
\SourceLeaf{408}\par\indent Die zur Ekliptik senkrechte Lage der Bahn ergab sich zwanglos aus den Beobachtungen, wird aber wohl durch weitere Beobachtungen noch etwas geändert werden. Erst dann wird man sehen, ob die Bewegung in Bezug auf die Ekliptik rechtläufig oder rückläufig ist. Um die Beobachtungen zu erleichtern, habe ich aus diesen Elementen die folgende Ephemeride berechnet. Die Zeit ist $14^{\mathrm u}20^{\mathrm m}$ in Göttingen. Die Spalte für die Helligkeit setzt geborgtes Licht voraus, $=1$ im Abstand 1 von Erde und Sonne.

\begin{center}\fontsize{10}{12}\selectfont\setlength{\tabcolsep}{3pt}
\begin{tabular}{ll|ll|lll|l}
\multicolumn{2}{c|}{1818} & \multicolumn{2}{c|}{Rekt.} & \multicolumn{3}{c|}{Dekl.} & Helligkeit \\
\hline
April & 6  & $298^\circ$ & $58'$ & $12^\circ$ & $57'$ & N. & 0.396 \\
      & 10 & 297 & 50 & 11 & 2  &    & 0.418 \\
      & 14 & 296 & 27 & 8  & 54 &    & 0.443 \\
      & 18 & 294 & 47 & 6  & 30 &    & 0.471 \\
      & 22 & 292 & 44 & 3  & 48 &    & 0.503 \\
      & 26 & 290 & 17 & 0  & 43 &    & 0.535 \\
      & 30 & 287 & 21 & 2  & 46 & S. & 0.568 \\
Mai   & 4  & 283 & 53 & 6  & 41 &    & 0.596 \\
      & 8  & 279 & 47 & 10 & 58 &    & 0.616 \\
      & 12 & 275 & 4  & 15 & 33 &    & 0.621 \\
\end{tabular}
\end{center}

\par\indent Wegen des Mondscheins wird man die Beobachtungen wohl nicht länger fortsetzen können. Wenn er aufhört, wird der Komet zumindest für uns zu weit nach Süden vorgerückt sein.

\par\indent Ich selbst habe den Kometen bisher erst einmal beobachtet und mit einem Stern 9\textsuperscript{ter} Größe verglichen, der in der Hist. cél. nicht vorkommt. Die Bestimmung seines Orts behalte ich mir für später vor. Er dürfte ungefähr bei $300^\circ\,3'$ Rektaszension und $15^\circ\,9'$ Deklination liegen.

\begin{center}\fontsize{9}{11}\selectfont\setlength{\tabcolsep}{3pt}\begin{tabular}{lrrrll}
1. April. & $15^{\mathrm u}$ & $45^{\mathrm m}$ & $28^{\mathrm s}$ & Sternzeit, & \shortstack[l]{der Komet folgt dem Stern um\\$14''98$ in Zeit}\\
& $15$ & $44$ & $56$ & — & \shortstack[l]{der Komet steht $4'\,22''$\\weiter südlich als der Stern.}\\
\end{tabular}\end{center}

\par\indent Die Deklination dürfte erheblich weniger genau sein als die Rektaszension.

\Needspace{8\baselineskip}\par\indent [Handschrift:] Beobachtungen dieses Sterns
\begin{center}\fontsize{10}{12}\selectfont\begin{tabular}{lrl}
1819 August & 27 & $20^{\mathrm u}\;0^{\mathrm m}\;11^{\mathrm s}49$\\
Sept. & 4 & $11.48$\\
& 7 & $11.38$\\
& 24 & $11.30$\\
\end{tabular}\end{center}

